% Part: first-order-logic
% Chapter: syntax-and-semantics
% Section: free-vars-sentences

\documentclass[../../../include/open-logic-section]{subfiles}

\begin{document}

\olfileid{fol}{syn}{fvs}

\olsection{মুক্ত \printtoken{P}{variable} ও \printtoken{P}{sentence}}

\begin{defn}[!!a{variable}-এর মুক্ত সংঘটন]
\ollabel{defn:free-occ}
!!a{formula}-এ !!a{variable}-এর \emph{মুক্ত} সংঘটনগুলি নিচের নিয়মে
আরোহীভাবে সংজ্ঞায়িত:
\begin{enumerate}
\item \indcase*{!A}{$!A$ পরমাণু}{$\indfrm$-এ !!{variable}-এর সব
  সংঘটন মুক্ত।}

\tagitem{prvNot}{\indcase{!A}{\lnot !B}{$\indfrm$-এ !!{variable}-এর
  মুক্ত সংঘটনগুলি ঠিক $!B$-এর মুক্ত সংঘটনগুলিই।}}{}

\item \indcase{!A}{(!B \ast !C)}{$\indfrm$-এ !!{variable}-এর মুক্ত
  সংঘটনগুলি হলো $!B$-এর মুক্ত সংঘটন এবং $!C$-এর মুক্ত সংঘটন একত্রে।}

\tagitem{prvAll}{\indcase{!A}{\lforall[x][!B]}{$\indfrm$-এর মুক্ত
  !!{variable}-সংঘটনগুলি হলো $!B$-এর সব মুক্ত সংঘটন থেকে $x$-এর
  সংঘটনগুলি বাদ দিলে যা থাকে।}}{}

\tagitem{prvEx}{\indcase{!A}{\lexists[x][!B]}{$\indfrm$-এর মুক্ত
  !!{variable}-সংঘটনগুলি হলো $!B$-এর সব মুক্ত সংঘটন থেকে $x$-এর
  সংঘটনগুলি বাদ দিলে যা থাকে।}}{}
\end{enumerate}
\end{defn}

\begin{defn}[বদ্ধ চলরাশি]
কোনো সূত্র~$!A$-এ !!a{variable}-এর যে সংঘটন মুক্ত নয়, সেটি
\emph{বদ্ধ}।
\end{defn}

\begin{prob}
\olref[fol][syn][fvs]{defn:free-occ}-এর অনুরূপভাবে বদ্ধ চলরাশির
সংঘটনগুলির একটি আরোহী সংজ্ঞা দিন।
\end{prob}

\begin{defn}[পরিসর]
\iftag{prvAll}{$\lforall[x][!B]$ যদি সূত্র~$!A$-এ কোনো উপসূত্রের
  সংঘটন হয়, তাহলে $!A$-এ $!B$-এর অনুরূপ সংঘটনটিকে $\lforall[x]$-এর
  অনুরূপ সংঘটনের \emph{পরিসর} বলা হয়। \iftag{prvEx}{$\lexists[x]$-এর
  ক্ষেত্রেও একই কথা।}{}}{যদি $\lexists[x][!B]$ সূত্র~$!A$-এ কোনো
  উপসূত্রের সংঘটন হয়, তাহলে $!A$-এ $!B$-এর অনুরূপ সংঘটনটিকে
  $\lexists[x]$-এর অনুরূপ সংঘটনের \emph{পরিসর} বলা হয়।}

$!B$ যদি~$!A$-এ কোনো পরিমাণসূচক-সংঘটন
\iftag{prvAll}{$\lforall[x]$\iftag{prvEx}{ বা
    $\lexists[x]$}{}}{$\lexists[x]$}-এর পরিসর হয়, তাহলে $!B$-তে
$x$-এর মুক্ত সংঘটনগুলি \iftag{prvAll}{$\lforall[x][!B]$\iftag{prvEx}{
ও $\lexists[x][!B]$}{}}{$\lexists[x][!B]$}-তে বদ্ধ। আমরা বলি, এই
সংঘটনগুলি উল্লিখিত পরিমাণসূচক-সংঘটনটির \emph{দ্বারা বদ্ধ}।
\end{defn}

\begin{ex}
নিচের সূত্রটি বিবেচনা করুন:
\[
\lexists[\Obj v_0][\underbrace{\Atom{\Obj A^2_0}{\Obj v_0,\Obj v_1}}_{!B}]
\]
$!B$ হলো $\lexists[\Obj v_0]$-এর পরিসর। পরিমাণসূচকটি $!B$-তে
$\Obj v_0$-এর সংঘটনকে বদ্ধ করে, কিন্তু $\Obj v_1$-এর সংঘটনকে বদ্ধ
করে না। তাই এই ক্ষেত্রে $\Obj v_1$ একটি মুক্ত চলরাশি।


এবার দেখি, আরও জটিল কোনো !!{formula}~$!A$-এ এটি কীভাবে কাজ করে:
\[
\lforall[\Obj v_0][\underbrace{(\Atom{\Obj A^1_0}{\Obj v_0} \lif
    \Atom{\Obj A^2_0}{\Obj v_0, \Obj v_1})}_{!B}] \lif \lexists[\Obj
  v_1][\underbrace{(\Atom{\Obj A^2_1}{\Obj v_0, \Obj v_1} \lor \lforall[\Obj v_0][\overbrace{\lnot \Atom{\Obj A^1_1}{\Obj v_0}}^{!D}])}_{!C}]
\]
$!B$ হলো প্রথম $\lforall[\Obj v_0]$-এর পরিসর, $!C$ হলো
$\lexists[\Obj v_1]$-এর পরিসর এবং $!D$ হলো দ্বিতীয়
$\lforall[\Obj v_0]$-এর পরিসর। প্রথম $\lforall[\Obj v_0]$,
$!B$-তে $\Obj v_0$-এর সংঘটনগুলিকে বদ্ধ করে; $\lexists[\Obj v_1]$,
$!C$-তে $\Obj v_1$-এর সংঘটনটিকে বদ্ধ করে; এবং দ্বিতীয়
$\lforall[\Obj v_0]$, $!D$-তে $\Obj v_0$-এর সংঘটনটিকে বদ্ধ করে।
$\Obj v_1$-এর প্রথম সংঘটন এবং $\Obj v_0$-এর চতুর্থ সংঘটন~$!A$-তে
মুক্ত। $\Obj v_0$-এর শেষ সংঘটনটি $!D$-তে মুক্ত, কিন্তু $!C$ ও~$!A$-তে
বদ্ধ।
\end{ex}

\begin{defn}[বাক্য]
!!^a{formula}~$!A$ তখনই \article{sentence} \emph{!!{sentence}}, যখন
তাতে !!{variable}s-এর কোনো মুক্ত সংঘটন নেই।
\end{defn}

% add examples!

\end{document}
